% Chapter fragment: load amsmath and graphicx in the main report.
% Figure paths support compilation from the repository root or problem02/.
\chapter{Deterministic and Stochastic Nonlinear Modeling}
\label{chap:problem2}

\section{Deterministic Nonlinear Model}

The states are the water masses $x=[m_1,m_2,m_3,m_4]^T$, the
manipulated inputs are $u=[F_1,F_2]^T$, and the unmeasured disturbances
are $d=[F_3,F_4]^T$. Flows are volumetric flow rates in
$\mathrm{cm^3/s}$, masses are in grams, and tank levels are in centimetres.
For constant density $\rho$ and tank cross-sectional areas $A_i$,
the levels and outlet flows satisfy
\begin{equation}
 h_i=\frac{m_i}{\rho A_i},\qquad
 q_i=a_i\sqrt{2g h_i},\qquad i=1,\ldots,4,
 \label{eq:p2-torricelli}
\end{equation}
where $a_i$ is the outlet area and $g$ is gravitational acceleration.
This model assumes incompressible water, constant tank areas, and ideal
gravity-driven discharge according to Torricelli's law.

Pump 1 sends the fraction $\gamma_1$ to tank 1 and the remainder to
tank 4. Pump 2 sends the fraction $\gamma_2$ to tank 2 and the remainder
to tank 3. Tanks 3 and 4 drain into tanks 1 and 2, respectively.
The mass balances are therefore
\begin{align}
 \dot m_1&=\rho(-q_1+q_3+\gamma_1F_1),\\
 \dot m_2&=\rho(-q_2+q_4+\gamma_2F_2),\\
 \dot m_3&=\rho(-q_3+(1-\gamma_2)F_2+F_3),\\
 \dot m_4&=\rho(-q_4+(1-\gamma_1)F_1+F_4).
 \label{eq:p2-balances}
\end{align}
Writing $q(x,p)=[q_1,q_2,q_3,q_4]^T$ gives the required model
\begin{equation}
 \dot x=f(x,u,d,p)=\rho\bigl(Qq(x,p)+Uu+Dd\bigr),
 \label{eq:p2-drift}
\end{equation}
with routing matrices
\begin{equation}
 Q=\begin{bmatrix}-1&0&1&0\\0&-1&0&1\\0&0&-1&0\\0&0&0&-1\end{bmatrix},
 \quad U=\begin{bmatrix}\gamma_1&0\\0&\gamma_2\\0&1-\gamma_2\\1-\gamma_1&0\end{bmatrix},
 \quad D=\begin{bmatrix}0&0\\0&0\\1&0\\0&1\end{bmatrix}.
\end{equation}
Internal flows cancel when the balances are summed:
$\sum_i\dot m_i=\rho(F_1+F_2+F_3+F_4-q_1-q_2)$.
Only the square-root discharge law introduces nonlinearity.

The sensors measure the lower-tank levels, which are also the controlled
outputs. In the deterministic model,
\begin{equation}
 y=g(x,p)=Cx,\qquad z=h(x,p)=Cx,\qquad
 C=\begin{bmatrix}(\rho A_1)^{-1}&0&0&0\\0&(\rho A_2)^{-1}&0&0\end{bmatrix}.
 \label{eq:p2-output}
\end{equation}
Here $h(x,p)$ denotes the output map, while $h_i$ denotes a tank level.
The physical parameters used in all experiments are
\begin{equation}
\begin{gathered}
 A_i=380.1327\ \mathrm{cm^2},\quad a_i=1.2272\ \mathrm{cm^2},\quad
 g=981\ \mathrm{cm/s^2},\\
 \rho=1\ \mathrm{g/cm^3},\qquad \gamma_1=0.6,\qquad\gamma_2=0.7.
\end{gathered}
\end{equation}

\section{Stochastic Nonlinear Model}

For piecewise constant disturbances, the dynamics retain the form
\eqref{eq:p2-drift}, but
\begin{equation}
 d(t)=d_k=\bar d+\eta_k,\qquad t_k\leq t<t_{k+1},\qquad
 \eta_k\sim\mathcal N(0,R_{dd}).
\end{equation}
The disturbance samples are independent across time and between the two
upper tanks. Conditional on a sampled disturbance sequence, the state
trajectory is obtained by solving an ordinary differential equation.
In the experiment,
\begin{equation}
 \bar d=\begin{bmatrix}20\\20\end{bmatrix}\ \mathrm{cm^3/s},\qquad
 R_{dd}=\begin{bmatrix}25&0\\0&25\end{bmatrix}(\mathrm{cm^3/s})^2,
\end{equation}
and each realization is held for $10\ \mathrm{s}$. A Gaussian distribution
allows negative flows; it approximates fluctuations around the positive
mean here, rather than imposing a strictly nonnegative inflow model.

The sampled sensor model and true controlled outputs are
\begin{equation}
 y_k=Cx(t_k)+v_k,\qquad v_k\sim\mathcal N(0,R_{vv}),\qquad
 z(t)=Cx(t),\qquad R_{vv}=0.25I_2\ \mathrm{cm^2}.
 \label{eq:p2-noisy-sensor}
\end{equation}
Measurement errors are independent across samples, between sensors, and
of the process disturbances. They affect observations but do not enter
the tank mass balances. Thus $z$ and $y$ must be distinguished even though
they describe the same levels.

\section{Stochastic Nonlinear Model (SDE)}

An alternative models accumulated disturbance volume as a mean inflow
plus Brownian fluctuations. Let $W=[W_3,W_4]^T$ consist of two independent
standard Wiener processes. The implemented It\^o SDE is
\begin{equation}
 dx=f(x,u,\bar d,p)\,dt+\Sigma\,dW,\qquad
 \Sigma=\rho D\begin{bmatrix}s_3&0\\0&s_4\end{bmatrix},\qquad
 s_3=s_4=5\ \mathrm{cm^3/s^{1/2}}.
 \label{eq:p2-sde}
\end{equation}
The mean disturbances remain in the drift. Over an interval $\Delta t$,
the random mass increment has covariance
\begin{equation}
 \operatorname{Cov}(\Sigma\Delta W)=\Sigma\Sigma^T\Delta t,\qquad
 \Sigma\Sigma^T=\rho^2\operatorname{diag}(0,0,s_3^2,s_4^2).
\end{equation}
Process noise enters tanks 3 and 4 directly and propagates to the lower
tanks through their outlet flows. The measurement model remains
\eqref{eq:p2-noisy-sensor}; measurement noise is sampled Gaussian noise,
whereas $W$ drives the continuous-time state dynamics.

Brownian volume increments represent rapid unresolved inflow fluctuations
and give a simple model with independent increments and constant diffusion.
They do not define an ordinary instantaneous flow $F_i(t)$: formally the
fluctuating flow is white noise. If inflow fluctuations have appreciable
temporal correlation, a mean-reverting disturbance state would be more
appropriate. Also, the value 5 in this diffusion model has different units
from the standard deviation 5 in the held-flow model. A held fluctuation
with variance $25$ contributes volume variance $25\Delta t^2$ over one
holding interval, while this SDE contributes $25\Delta t$ (with the
respective units). The runs illustrate different disturbance assumptions
rather than matched process-noise intensities.

\section{Simulation}

\subsection{Experimental setup and numerical method}

The three open-loop experiments start with empty tanks, $x(0)=0$, and run
until $1200\ \mathrm{s}$. The pump inputs are held constant between samples:
\begin{equation}
 F_1(t)=\begin{cases}300,&t<t_a,\\350,&t\geq t_a,\end{cases}
 \qquad
 F_2(t)=\begin{cases}300,&t<t_b,\\250,&t\geq t_b,\end{cases}
 \quad [\mathrm{cm^3/s}].
\end{equation}
The deterministic and held-disturbance experiments use a $10\ \mathrm{s}$
grid and MATLAB's \texttt{ode15s}, restarted on each interval to preserve
the input discontinuities. Their step times are $t_a=390\ \mathrm{s}$ and
$t_b=790\ \mathrm{s}$. The SDE uses a $0.1\ \mathrm{s}$ grid, giving
$t_a=399.9\ \mathrm{s}$ and $t_b=799.9\ \mathrm{s}$. These differences
follow from the sample-index-based step definition in
\texttt{problem02Experiments.m}.

The SDE is integrated with the projected Euler--Maruyama update
\begin{equation}
 x_{k+1}=\max\left\{0,\ x_k+f(x_k,u_k,\bar d,p)\Delta t
                    +\Sigma\Delta W_k\right\},\qquad
 \Delta W_k\sim\mathcal N(0,\Delta t I_2),
\end{equation}
where the maximum acts componentwise. Constant diffusion makes the
Milstein correction zero. The projection prevents negative masses near
empty tanks; it is a numerical boundary treatment and modifies the
unconstrained SDE at that boundary. The ODE drift evaluates the square root
using nonnegative trial masses. The deterministic run uses $d=\bar d$ and
zero measurement noise. Both stochastic runs use measurement standard
deviation $0.5\ \mathrm{cm}$ and MATLAB's \texttt{twister} generator with
seed 100. Their sampling grids and random-array sizes differ, so the seed
does not make them share the same noise realization.

\subsection{Open-loop results}

\begingroup
\newcommand{\pTwoFigure}[1]{%
 \IfFileExists{results/problem02/#1.pdf}{%
  \includegraphics[width=\linewidth]{problem02/results/#1.pdf}}{%
  \includegraphics[width=\linewidth]{results/#1.pdf}}}
\begin{figure}[htbp]
 \centering
 \pTwoFigure{open_loop_01}
 \caption{Deterministic open-loop simulation. Pump inputs and constant
 disturbance flows are shown above the lower-tank outputs. With zero
 measurement noise, $y_i=z_i$.}
 \label{fig:p2-deterministic}
\end{figure}
\begin{figure}[htbp]
 \centering
 \pTwoFigure{open_loop_02}
 \caption{Open-loop simulation with Gaussian disturbance flows held for
 $10\ \mathrm{s}$ and independent measurement noise. The curves $z_i$
 show true levels; the points $y_i$ are noisy observations.}
 \label{fig:p2-held}
\end{figure}
\begin{figure}[htbp]
 \centering
 \pTwoFigure{open_loop_03}
 \caption{SDE open-loop simulation with $0.1\ \mathrm{s}$ integration steps.
 The middle panel shows upper-tank levels $h_3,h_4$, which receive the
 Brownian process noise directly. The bottom panel shows measured and
 true lower-tank levels.}
 \label{fig:p2-sde}
\end{figure}
\endgroup

Figures~\ref{fig:p2-deterministic}--\ref{fig:p2-sde} show filling from the
empty initial state and the responses to the two pump steps. Increasing
$F_1$ raises tank 1 through its direct inflow and tank 2 through tank 4.
Decreasing $F_2$ lowers tank 2 directly and tank 1 through tank 3.
The indirect paths include an additional tank and produce a slower response.
In the stochastic runs, the lower tanks smooth the upper-tank disturbances,
while measurement noise produces scatter around the true outputs.

Table~\ref{tab:p2-results} reports the final true levels and empirical
root-mean-square measurement errors, computed as
$\sqrt{N^{-1}\sum_{k=1}^N(y_{i,k}-z_{i,k})^2}$ from
\texttt{results/open\_loop.mat}.
\begin{table}[htbp]
 \centering
 \caption{Final levels at $1200\ \mathrm{s}$ and measurement-error RMS,
 all in centimetres. Each stochastic result is a single realization.}
 \label{tab:p2-results}
 \begin{tabular}{lrrrrrr}
  \hline
  Model & $h_1$ & $h_2$ & $h_3$ & $h_4$ & RMS$_1$ & RMS$_2$\\
  \hline
  Deterministic & 31.509 & 38.056 & 3.054 & 8.664 & 0.000 & 0.000\\
  Held disturbances & 31.498 & 37.943 & 3.055 & 8.357 & 0.471 & 0.549\\
  SDE & 31.469 & 38.074 & 3.001 & 8.642 & 0.502 & 0.497\\
  \hline
 \end{tabular}
\end{table}

As a deterministic reference, setting the balances to zero gives
\begin{align}
 q_3^*&=(1-\gamma_2)F_2+\bar F_3,&
 q_4^*&=(1-\gamma_1)F_1+\bar F_4,\\
 q_1^*&=\gamma_1F_1+q_3^*,&
 q_2^*&=\gamma_2F_2+q_4^*,\\
 h_i^*&=\frac{(q_i^*/a_i)^2}{2g}.&&
\end{align}
For the final inputs $u=[350,250]^T$ and $\bar d=[20,20]^T$, this yields
$h^*=[31.483,37.980,3.054,8.664]^T\ \mathrm{cm}$.
The deterministic final lower-tank levels remain slightly above this
equilibrium because the response to the second step has not completely
settled. The stochastic final levels are close to this reference, and
their measurement RMS values are consistent with the specified
$0.5\ \mathrm{cm}$ standard deviation. These single trajectories illustrate
model behaviour; they do not establish ensemble means, confidence
intervals, or numerical convergence of the SDE discretization.


testing
